\section{Discussion}
\label{sec:discussion}

\subsection{Interpretation of Results}

Our results provide converging evidence that self-critique text induces genuine geometric restructuring of the residual stream, beyond what surface-level text variation produces. Five independent metrics---cosine similarity, $L_2$ displacement, CKA, drift consistency, and linear probes---all indicate that critique activations diverge from base activations more than paraphrase activations do. The effect is statistically robust (29/33 layers significant after Bonferroni correction) and geometrically structured (consistent drift direction across samples).

\para{The critique subspace.}
The finding that critique drift is more directionally consistent than paraphrase drift (0.596 vs.\ 0.479 in layers 10--30) suggests the existence of a systematic ``critique subspace'' in the residual stream. This is consistent with \citet{zhu2025emergence}, who identify separable self-reflection vectors in hidden states. Our contribution is showing that this directional structure emerges even in a base (non-instruction-tuned) model processing externally provided critique text, suggesting that the capacity to distinguish evaluative from non-evaluative content may be a general property of pretrained language models.

\para{Layer specificity.}
The concentration of maximal drift in middle layers (12--16, approximately 40--50\% of network depth) aligns with \citet{polo2026emergent}, who find that middle layers initiate geometric restructuring during reasoning. This suggests a common computational pattern: middle layers perform the heavy lifting of restructuring representations in response to new information, while early layers encode surface features and later layers consolidate task-relevant representations.

\subsection{Limitations}

\para{Text content confound.}
The critique and paraphrase conditions differ in both structure \emph{and} content. Critique text explicitly discusses correctness, which inflates linear probe accuracy and may contribute to the larger drift magnitude. A stronger control would use critique-structured text that evaluates irrelevant aspects of the answer (\eg style rather than correctness), isolating the effect of evaluative structure from correctness content.

\para{Externally generated stimuli.}
Critique text was generated by GPT-4.1, not by \pythia itself. Our experiment therefore measures how \pythia \emph{processes} critique text, not how it \emph{generates} self-critique. A complementary study should have the model generate its own critiques, though this is challenging with base (non-instruction-tuned) models.

\para{Single model.}
All results are from \pythia only. Generalization to instruction-tuned models, reasoning-specialized models (\eg DeepSeek-R1), and different model scales remains to be established. Instruction-tuned models may show amplified or qualitatively different critique processing dynamics.

\para{Last-token extraction.}
We extract activations only at the last token position. While this captures the model's integrated representation of the full input, token-level analysis across the sequence could reveal richer dynamics---for instance, whether critique processing is localized to specific token positions.

\para{Correlational evidence.}
We measure correlational differences in activations but do not perform causal interventions. Activation steering along the identified critique drift direction would test whether this direction functionally alters model behavior, strengthening the evidence for a genuine ``critique subspace.''

\subsection{Broader Implications}

\para{For AI systems design.}
Our results support the use of self-reflection prompts in architectures that rely on genuine self-evaluation. The finding that critique produces measurably different internal states---not just different outputs---suggests that reflection-based systems like Reflexion~\citep{shinn2023reflexion} may benefit from real representational change rather than mere output re-sampling.

\para{For mechanistic interpretability.}
The identified critique subspace is a candidate for causal investigation via activation steering~\citep{zou2023representation}. If steering along the critique direction can induce or suppress self-evaluative behavior, this would connect geometric observations to functional mechanisms.

\para{For AI safety.}
Understanding whether recursive self-improvement systems produce genuine cognitive change is critical for safety evaluation. Our finding that critique creates genuinely different representations (not just re-sampling) is relevant for assessing whether self-improvement loops can compound changes in model behavior over multiple iterations.
